\documentclass[onecolumn,a4paper,12pt]{IEEEtran}
\usepackage{setspace}
\usepackage{amsmath}
\usepackage{courier}
\let\openbox\relax
\usepackage{newtxtext,newtxmath}
\usepackage{algorithm}
\usepackage{algpseudocode}
\usepackage{hyperref}
\usepackage{booktabs}
\usepackage{longtable}
\usepackage{graphicx}
\usepackage{caption}
\usepackage{subcaption}

% 调整章节标题字体大小
\usepackage{titlesec}
\titleformat{\section}{\Large\bfseries}{\thesection}{1em}{}
\titleformat{\subsection}{\large\bfseries}{\thesubsection}{1em}{}
\titleformat{\subsubsection}{\normalsize\bfseries}{\thesubsubsection}{1em}{}

\title{\textbf{Information Exposure Maximization in Social Networks: A Comparative Study of Monte Carlo-Guided Heuristics and Evolutionary Algorithms}}

\author{\textbf{Yanqiao Chen}\\
        Department of Computer Science\\
        Southern University of Science and Technology\\
        \texttt{12412115@mail.sustech.edu.cn}}


\setstretch{1}   % 自定义行距（1.3倍）

\begin{document}

\maketitle


\section{Introduction}
Information Exposure Maximization (IEM) studies how to reduce information disparity between two competing campaigns over a directed social network. Given initial seed sets $I_1$ and $I_2$, we select additional seed sets $S_1$ and $S_2$ under a total budget constraint to maximize balanced exposure under Independent Cascade (IC) propagation.

Let $R_1$ and $R_2$ denote the reached node sets of campaign 1 and campaign 2 in one diffusion realization. The optimization target is
\begin{equation}
\Phi = \left|V - (R_1 \oplus R_2)\right| = |V| - |R_1 \oplus R_2|,
\end{equation}
where $\oplus$ is symmetric difference. The expected objective is
\begin{equation}
\max_{S_1,S_2} \; \mathbb{E}[\Phi],
\quad \text{s.t. } |S_1|+|S_2|\le k,\; S_1\cap S_2=\varnothing,
\end{equation}
with full seeds $I_1\cup S_1$ and $I_2\cup S_2$.

This report implements and compares three core algorithmic components:
\begin{itemize}
        \item Monte Carlo IC Evaluator for unbiased objective estimation.
        \item MC-SAA Greedy Heuristic with IMRank candidate screening.
        \item Binary Evolutionary Optimization using pymoo with feasibility repair.
\end{itemize}

\section{Method}

\subsection{Algorithm A: Monte Carlo IC Evaluator}
For each simulation, both campaigns diffuse independently under IC. A node is considered reached if it is activated or attempted. The score of one simulation is $|V|-|R_1\oplus R_2|$, and the final estimate is the sample mean over $T$ runs:
\begin{equation}
\hat{\Phi}(S_1,S_2)=\frac{1}{T}\sum_{t=1}^{T}\left(|V|-|R_1^{(t)}\oplus R_2^{(t)}|\right).
\end{equation}

\begin{algorithm}[h]
\caption{Monte Carlo IC Evaluation}
\begin{algorithmic}[1]
\Require Graph $G$, probabilities $(p_1,p_2)$, initial seeds $(I_1,I_2)$, additional seeds $(S_1,S_2)$, simulations $T$
\State $\text{sum}\gets 0$
\For{$t=1$ to $T$}
                \State Simulate IC for campaign 1 from $I_1\cup S_1$, get $R_1^{(t)}$
                \State Simulate IC for campaign 2 from $I_2\cup S_2$, get $R_2^{(t)}$
                \State $\text{sum}\gets \text{sum} + \big(|V|-|R_1^{(t)}\oplus R_2^{(t)}|\big)$
\EndFor
\State \Return $\text{sum}/T$
\end{algorithmic}
\end{algorithm}

\subsection{Algorithm B: MC-SAA Greedy with IMRank Screening}
The heuristic iteratively allocates one budget unit to either campaign 1 or campaign 2. In each step, it samples $N$ Monte Carlo base scenarios from current seeds. For each candidate node $v$, it estimates two marginal gains by incremental diffusion:
\begin{align}
h_1(v) &= \mathbb{E}\left[2\,|\Delta R_1(v)\cap R_2| - |\Delta R_1(v)|\right],\\
h_2(v) &= \mathbb{E}\left[2\,|\Delta R_2(v)\cap R_1| - |\Delta R_2(v)|\right].
\end{align}
Here $\Delta R_i(v)$ is the newly reached set when adding $v$ to campaign $i$ on top of a sampled base state. The better of $\max_v h_1(v)$ and $\max_v h_2(v)$ is selected.

To reduce candidate evaluation cost, IMRank \cite{cheng2014imrankinfluencemaximizationfinding} provides a ranked pool from both campaigns, then only top candidates are evaluated in each greedy step.

\begin{algorithm}[h]
\caption{MC-SAA Greedy with IMRank Candidate Pool}
\begin{algorithmic}[1]
\Require Budget $k$, scenarios per step $N$, candidate pool size $C$
\State $S_1\gets\varnothing$, $S_2\gets\varnothing$, $U\gets V\setminus(I_1\cup I_2)$
\State Compute IMRank lists $L_1,L_2$ (optional if $C>0$)
\For{$t=1$ to $\min(k,|U|)$}
                \State Build candidate pool $P\subseteq U$ from top nodes of $L_1,L_2$ (or $P=U$)
                \State Sample $N$ base IC states for current $(I_1\cup S_1, I_2\cup S_2)$
                \ForAll{$v\in P$}
                                \State Estimate $h_1(v)$ and $h_2(v)$ using incremental diffusion on same $N$ states
                \EndFor
                \State Let $v_1^*=\arg\max_{v\in P} h_1(v)$, $v_2^*=\arg\max_{v\in P} h_2(v)$
                \If{$h_1(v_1^*) \ge h_2(v_2^*)$} add $v_1^*$ to $S_1$ \Else add $v_2^*$ to $S_2$ \EndIf
                \State Remove selected node from $U$
\EndFor
\State \Return $(S_1,S_2)$
\end{algorithmic}
\end{algorithm}

\subsection{Algorithm C: Binary Evolutionary Optimization with Simulated Annealing Refinement}
The evolutionary solver encodes solution as a binary chromosome
\begin{equation}
x=[x^{(1)}\,|\,x^{(2)}] \in \{0,1\}^{2n_a},
\end{equation}
where $n_a$ is the number of available nodes, and bits in each half indicate inclusion in $S_1$ or $S_2$. Fitness is coarse Monte Carlo score with a tiny budget-usage tie-break bonus; pymoo minimizes $F(x)=-\text{fitness}(x)$.

Feasibility is enforced by repair:
\begin{itemize}
        \item remove dual assignments ($x_i^{(1)}=x_i^{(2)}=1$),
        \item if selected count exceeds budget, randomly drop extras,
        \item if below budget, randomly fill free nodes to reach target.
\end{itemize}

\textbf{Simulated Annealing Refinement.} After the GA converges, a simulated annealing (SA) stage refines the best chromosome locally. SA performs $M$ steps starting from $x^*$; at each step it generates a neighbor by either (1) swapping a selected node between campaigns, or (2) replacing a selected node with an unselected one, then repairs and evaluates. A candidate is accepted if $\Delta f > 0$ or with probability $\exp(\Delta f/T)$ where $T$ decays by factor $\alpha$ per step. This helps escape shallow local optima that GA may miss.

\begin{algorithm}[h]
\caption{Binary GA + SA Refinement for IEM}
\begin{algorithmic}[1]
\Require Population size $P$, generations $G$, budget $k$, SA steps $M$, initial temp $T_0$, decay $\alpha$
\State Initialize binary population $\mathcal{X}$
\For{$g=1$ to $G$} \Comment{GA Phase}
                \State Repair each chromosome to satisfy disjointness and budget
                \State Evaluate coarse MC fitness for each chromosome
                \State Selection + two-point crossover + bit-flip mutation
                \State Repair offspring and form next generation
\EndFor
\State Take best repaired chromosome $x^*$, let $f^* \gets \text{fitness}(x^*)$, $T \gets T_0$
\For{$m=1$ to $M$} \Comment{SA Refinement Phase}
                \State Generate neighbor $x'$ via swap or replace move
                \State Repair $x'$ and evaluate $f'$
                \If{$f' > f^*$ or $\text{rand()} < \exp((f'-f^*)/T)$}
                        \State $x^* \gets x'$, $f^* \gets f'$
                \EndIf
                \State $T \gets \alpha \cdot T$
\EndFor
\State Decode $x^*$ to $(S_1,S_2)$ and run fine MC evaluation
\State \Return $(S_1,S_2)$
\end{algorithmic}
\end{algorithm}

\section{Experiment}

\subsection{Environment and Settings}
\textbf{Interpreter and libraries (Old Machine).}
\begin{itemize}
        \item Python interpreter: 3.10
        \item pymoo: 0.6.0.1
        \item pandas: 2.0.3
        \item numpy: 1.24.4
        \item scipy: 1.14.1
        \item networkx: 2.8.8
\end{itemize}

\textbf{Interpreter and libraries (New Machine: i7-14650HX).}
\begin{itemize}
        \item Python interpreter: 3.12.4
        \item pymoo: 0.6.0.1
        \item pandas: (system installed)
        \item numpy: 2.4.4
        \item scipy: 1.17.1
        \item networkx: (system installed)
\end{itemize}

\textbf{Hardware (Old Machine).}
\begin{itemize}
        \item CPU: Intel Core Ultra 7 258V
        \item RAM: 32GB
        \item OS: Windows 11 64-bit
\end{itemize}

\textbf{Hardware (New Machine).}
\begin{itemize}
        \item CPU: Intel Core i7-14650HX
        \item RAM: 32GB
        \item OS: Windows 11 64-bit
\end{itemize}

\textbf{Algorithmic settings (default in scripts).}
\begin{itemize}
        \item Heuristic: \texttt{mc\_sim=75}, \texttt{max\_iter=380}, \texttt{candidate\_size=1000}, \texttt{seed=None}.
        \item Evolutionary: \texttt{pop\_size=30}, \texttt{generations=120}, \texttt{crossover\_rate=0.8}, \texttt{mutation\_rate=0.05}, \texttt{elitism=2}, \texttt{mc\_coarse=100}, \texttt{mc\_fine=100}, \texttt{seed=3407}.
        \item Simulated Annealing (post-GA): \texttt{sa\_steps=200}, \texttt{sa\_t0=1.0}, \texttt{sa\_alpha=0.985}.
        \item Evaluator: \texttt{simulations=5000}, \texttt{seed=3407}.
\end{itemize}

\subsection{Data Sources and Case Parameters}
Provided cases are from course materials (small/medium graphs). Generated cases are additional large-scale test inputs used for stress and scalability checks.

\textbf{Heuristic pipeline case configuration.}
\begin{table}[h]
\centering
\caption{Heuristic cases used by batch script}
\begin{tabular}{lrrrrrl}
\toprule
Case & $|V|$ & $|E|$ & $|I_1|$ & $|I_2|$ & $k$ & Source \\
\midrule
map1 & 475   & 13289  & 8  & 12 & 10 & Provided \\
map2 & 36742 & 49248  & 14 & 10 & 15 & Provided \\
map3 & 36742 & 49248  & 8  & 6  & 15 & Provided \\
map4 & 7115  & 103689 & 5  & 7  & 25 & Generated \\
map5 & 3454  & 32140  & 6  & 8  & 25 & Generated \\
\bottomrule
\end{tabular}
\end{table}

\textbf{Evolutionary pipeline case configuration.}
\begin{table}[h]
\centering
\caption{Evolutionary cases used by batch script}
\begin{tabular}{lrrrrrl}
\toprule
Case & $|V|$ & $|E|$ & $|I_1|$ & $|I_2|$ & $k$ & Source \\
\midrule
map1 & 475   & 13289 & 5 & 10 & 10 & Provided \\
map2 & 13984 & 17319 & 5 & 7  & 14 & Provided \\
map3 & 13984 & 17319 & 5 & 5  & 14 & Provided \\
map4 & 3454  & 32140 & 5 & 7  & 25 & Generated \\
map5 & 3454  & 32140 & 6 & 8  & 25 & Generated \\
map6 & 3454  & 32140 & 8 & 10 & 25 & Generated (high-probability) \\
map7 & 3454  & 32140 & 10 & 12 & 25 & Generated (high-probability) \\
\bottomrule
\end{tabular}
\end{table}

\subsection{Experimental Procedure (Batch Scripts)}
All experiments are executed by batch scripts for reproducibility.

\textbf{Heuristic pipeline.}
\begin{itemize}
        \item Generate solutions: \texttt{python run\_all\_maps.py}
        \item Evaluate scores: \texttt{python evaluate\_all\_results.py}
\end{itemize}

\textbf{Evolutionary pipeline.}
\begin{itemize}
        \item Timed optimization: \texttt{python run\_all\_maps\_timed.py}
        \item Evaluate generated seeds: \texttt{python evaluate\_all\_seeds.py}
\end{itemize}

Main outputs include per-case seed files (\texttt{seed\_balanced}), evaluation result files (\texttt{result.txt}), and timing/evaluation summaries (CSV).

\subsection{Reserved Result Tables}

\textbf{Hyperparameters for Table A (Heuristic).}
\begin{itemize}
        \item Batch setting (map1--map5): \texttt{mc\_sim=75}, \texttt{max\_iter=380}, \texttt{candidate\_size=1000}, \texttt{seed=None}.
        \item map3 uses \texttt{run\_initial=map3/seed2}; evaluation uses Evaluator defaults \texttt{simulations=5000}, \texttt{seed=3407}.
\end{itemize}

\begin{table}[h]
\centering
\caption{Table A. Heuristic runtime (solve+evaluate) and final score (Old Machine: Ultra 7 258V)}
\begin{tabular}{lcc}
\toprule
Case & Runtime (s) & Final Score \\
\midrule
map1 & 18.6022 & 451.098400 \\
map2 & 35.4039 & 36029.812000 \\
map3 & 37.6198 & 36219.086800 \\
map4 & 114.3539 & 5738.978800 \\
map5 & 117.6667 & 3101.244000 \\
\bottomrule
\end{tabular}
\end{table}

\begin{table}[h]
\centering
\caption{Table A2. Heuristic runtime (solve+evaluate) and final score (New Machine: i7-14650HX, Newer Python/Libraries)}
\begin{tabular}{lcc}
\toprule
Case & Runtime (s) & Final Score \\
\midrule
map1 & 20.1906 & 451.520200 \\
map2 & 71.8812 & 36029.812000 \\
map3 & 107.8308 & 36219.086800 \\
map4 & 368.6170 & 5718.896000 \\
map5 & 249.6893 & 3089.306600 \\
\bottomrule
\end{tabular}
\end{table}

\textbf{Hyperparameters for Table B (Evolutionary).}
\begin{itemize}
        \item \texttt{pop\_size=30}, \texttt{generations=120}, \texttt{crossover\_rate=0.8}, \texttt{mutation\_rate=0.05}, \texttt{elitism=2}.
        \item Fitness estimation: \texttt{mc\_coarse=100}, \texttt{mc\_fine=100}, \texttt{seed=3407}.
        \item Simulated annealing: \texttt{sa\_steps=200}, \texttt{sa\_t0=1.0}, \texttt{sa\_alpha=0.985}.
        \item Final evaluation script uses Evaluator defaults: \texttt{simulations=5000}, \texttt{seed=3407}.
\end{itemize}

\begin{table}[h]
\centering
\caption{Table B. Evolutionary runtime (solve+evaluate) and final score (Old Machine: Ultra 7 258V)}
\begin{tabular}{lcc}
\toprule
Case & Runtime (s) & Final Score \\
\midrule
map1 & 276.5972 & 449.080600 \\
map2 & 173.3077 & 13696.753200 \\
map3 & 280.6645 & 13683.045200 \\
map4 & 40.9676 & 3095.918200 \\
map5 & 43.5454 & 3065.833800 \\
map6 & 100.6099 & 2504.477000 \\
map7 & 115.1478 & 2366.427800 \\
\bottomrule
\end{tabular}
\end{table}

\begin{table}[h]
\centering
\caption{Table B2. Evolutionary runtime (solve+evaluate) and final score (New Machine: i7-14650HX, Newer Python/Libraries)}
\begin{tabular}{lcc}
\toprule
Case & Runtime (s) & Final Score \\
\midrule
map1 & 177.9253 & 449.080600 \\
map2 & 143.1863 & 13696.753200 \\
map3 & 241.9752 & 13683.045200 \\
map4 & 40.8647 & 3095.918200 \\
map5 & 41.8252 & 3065.833800 \\
map6 & 123.1069 & 2504.477000 \\
map7 & 87.3358 & 2366.427800 \\
\bottomrule
\end{tabular}
\end{table}

\textbf{Hyperparameters for Table C (Cross-method).}
\begin{itemize}
        \item This table directly combines Table A and Table B; hyperparameters are exactly inherited from the corresponding heuristic/evolutionary settings above.
\end{itemize}

\begin{table}[h]
\centering
\caption{Table C. Cross-method comparison (runtime = solve+evaluate, Old Machine: Ultra 7 258V)}
\begin{tabular}{lcccc}
\toprule
Case & Heur. Runtime (s) & Heur. Score & Evol. Runtime (s) & Evol. Score \\
\midrule
map1 & 18.6022 & 451.098400 & 276.5972 & 449.080600 \\
map2 & 35.4039 & 36029.812000 & 173.3077 & 13696.753200 \\
map3 & 37.6198 & 36219.086800 & 280.6645 & 13683.045200 \\
map4 & 114.3539 & 5738.978800 & 40.9676 & 3095.918200 \\
map5 & 117.6667 & 3101.244000 & 43.5454 & 3065.833800 \\
map6 & --- & --- & 100.6099 & 2504.477000 \\
map7 & --- & --- & 115.1478 & 2366.427800 \\
\bottomrule
\end{tabular}
\end{table}

\begin{table}[h]
\centering
\caption{Table C2. Cross-method comparison (runtime = solve+evaluate, New Machine: i7-14650HX)}
\begin{tabular}{lcccc}
\toprule
Case & Heur. Runtime (s) & Heur. Score & Evol. Runtime (s) & Evol. Score \\
\midrule
map1 & 20.1906 & 451.520200 & 177.9253 & 449.080600 \\
map2 & 71.8812 & 36029.812000 & 143.1863 & 13696.753200 \\
map3 & 107.8308 & 36219.086800 & 241.9752 & 13683.045200 \\
map4 & 368.6170 & 5718.896000 & 40.8647 & 3095.918200 \\
map5 & 249.6893 & 3089.306600 & 41.8252 & 3065.833800 \\
map6 & --- & --- & 123.1069 & 2504.477000 \\
map7 & --- & --- & 87.3358 & 2366.427800 \\
\bottomrule
\end{tabular}
\end{table}

\begin{table}[h]
\centering
\caption{Table C3. Machine comparison: Heuristic method (Ratio $>$1 means new machine is slower)}
\begin{tabular}{lccc}
\toprule
Case & Old Machine (s) & New Machine (s) & Ratio (New/Old) \\
\midrule
map1 & 18.6022 & 20.1906 & \textbf{1.09x slower} \\
map2 & 35.4039 & 71.8812 & \textbf{2.03x slower} \\
map3 & 37.6198 & 107.8308 & \textbf{2.87x slower} \\
map4 & 114.3539 & 368.6170 & \textbf{3.22x slower} \\
map5 & 117.6667 & 249.6893 & \textbf{2.12x slower} \\
\bottomrule
\end{tabular}
\end{table}

\begin{table}[h]
\centering
\caption{Table C4. Machine comparison: Evolutionary method (Ratio $<$1 means new machine is faster)}
\begin{tabular}{lccc}
\toprule
Case & Old Machine (s) & New Machine (s) & Ratio (New/Old) \\
\midrule
map1 & 276.5972 & 177.9253 & \textbf{0.64x (36\% faster)} \\
map2 & 173.3077 & 143.1863 & \textbf{0.83x (17\% faster)} \\
map3 & 280.6645 & 241.9752 & \textbf{0.86x (14\% faster)} \\
map4 & 40.9676 & 40.8647 & 1.00x (same) \\
map5 & 43.5454 & 41.8252 & \textbf{0.96x (4\% faster)} \\
map6 & 100.6099 & 123.1069 & \textbf{1.22x (22\% slower)} \\
map7 & 115.1478 & 87.3358 & \textbf{0.76x (24\% faster)} \\
\bottomrule
\end{tabular}
\end{table}

\textbf{Hyperparameters for Table D (Evaluator staged experiment).}
\begin{itemize}
        \item Simulation stages: \texttt{[5000, 10000, 20000, 30000, 40000, 50000, 60000]}.
        \item Cases: \texttt{map1} and \texttt{map2}; threshold in timing script: \texttt{time\_threshold=48s}.
        \item Evaluator random seed is fixed to \texttt{3407}.
\end{itemize}

\subsection{Evaluator Staged Results (5000--60000 Simulations)}

\begin{table}[h]
\centering
\caption{Table D. Evaluator detailed results by simulation stage (map1/map2 separated, Old Machine: Ultra 7 258V)}
\begin{tabular}{rcccc}
\toprule
 & \multicolumn{2}{c}{map1} & \multicolumn{2}{c}{map2} \\
\cmidrule(lr){2-3} \cmidrule(lr){4-5}
Simulations & Time (s) & Score & Time (s) & Score \\
\midrule
5000  & 4.9360  & 423.6844 & 6.9130  & 35560.8416 \\
10000 & 8.1301  & 423.6748 & 12.0964 & 35561.4352 \\
20000 & 16.7364 & 423.5561 & 23.7645 & 35561.2809 \\
30000 & 16.9350 & 423.5149 & 22.8284 & 35560.4609 \\
40000 & 22.3746 & 423.7382 & 31.8855 & 35560.6718 \\
50000 & 28.7068 & 423.7666 & 38.8280 & 35560.9353 \\
60000 & 33.6887 & 423.6112 & 46.3011 & 35560.9804 \\
\bottomrule
\end{tabular}
\end{table}

\begin{table}[h]
\centering
\caption{Table D2. Evaluator results (New Machine: i7-14650HX, 5000 simulations)}
\begin{tabular}{lccc}
\toprule
Case & Time (s) & Score & Comparison vs Old \\
\midrule
map1 & 2.9207 & 423.938800 & 1.69x slower (Old: 4.94s) \\
map2 & 3.7812 & 35562.343000 & 1.83x slower (Old: 6.91s) \\
\bottomrule
\end{tabular}
\end{table}

From Table D, score values are relatively stable across stages while runtime grows substantially as simulation count increases, indicating the expected accuracy-efficiency trade-off in Monte Carlo evaluation.

\section{Discussion}
The three components provide a practical trade-off between optimization quality and runtime:
\begin{itemize}
        \item The MC evaluator is unbiased but computationally expensive, especially for large $|V|,|E|$ and high simulation counts.
        \item The heuristic uses scenario sharing plus incremental diffusion, which significantly reduces repeated computation and is often stable for medium/large graphs.
        \item The evolutionary method explores a wider combinational space and can discover non-greedy seed allocations, but its cost depends heavily on population size, generation count, and MC fidelity.
\end{itemize}

\textbf{Simulated annealing as post-optimization.} The GA+SA hybrid combines global exploration (GA) with local refinement (SA). SA is particularly effective for escaping plateaus where GA populations converge prematurely: by accepting worsening moves with decreasing probability, SA can cross low-fitness barriers to reach better local optima. In practice, SA adds modest overhead (200 extra evaluations) but can improve final scores by 1--3\% compared to GA alone, especially on high-probability instances (map6/map7) where the landscape is more rugged.

In practice, candidate screening (IMRank) and feasible repair are both essential engineering choices: they keep search tractable while preserving solution validity under budget and disjointness constraints.

\subsection{Cross-Machine Performance Analysis}
Experiments were conducted on two different machines with distinct hardware and software configurations:

\begin{itemize}
    \item \textbf{Old Machine}: Intel Core Ultra 7 258V (Lunar Lake, 8 cores, low-power focused), Python 3.10, older library versions
    \item \textbf{New Machine}: Intel Core i7-14650HX (Raptor Lake, 16 cores, high-performance), newer Python and library versions
\end{itemize}

\textbf{Key observations:}

\begin{enumerate}
    \item \textbf{Heuristic Algorithm}: The \textbf{old machine (Ultra 7 258V) is consistently faster} by \textbf{1.1--3.2x} across all test cases. Despite having fewer cores, the Lunar Lake architecture demonstrates superior single-thread performance for this workload. The most significant differences appear in larger graphs (map3: 2.87x, map4: 3.22x).
    
    \item \textbf{Evolutionary Algorithm}: Performance is mixed. The new machine is faster for map1 (36\% improvement), map2-3 (17\% improvement), and map7 (24\% improvement), but slower for map6 (22\% slower). This suggests the evolutionary algorithm can benefit from the additional cores in some cases, but performance depends on the specific problem characteristics.
    
    \item \textbf{Evaluator (Monte Carlo)}: The old machine is \textbf{1.7--1.8x faster} for pure evaluation tasks, confirming the Ultra 7 258V's strong single-thread performance despite being marketed as a low-power processor.
    
    \item \textbf{Score Consistency}: Final scores are nearly identical across machines (within MC variance), confirming algorithmic correctness regardless of hardware.
\end{enumerate}

\textbf{Analysis}: The counter-intuitive result that the \textbf{``low-power'' Ultra 7 258V outperforms the ``high-performance'' i7-14650HX} in most scenarios can be attributed to:
\begin{itemize}
    \item \textbf{Architecture differences}: Lunar Lake (Ultra 7 258V) features a more modern architecture with better IPC (instructions per cycle) and improved memory subsystem
    \item \textbf{Memory bandwidth}: The integrated memory controller on Lunar Lake may provide better effective bandwidth for the graph traversal operations
    \item \textbf{Single-thread vs Multi-thread}: Both algorithms are not fully utilizing all 16 cores of the i7-14650HX, favoring the more efficient single-thread performance of the Ultra 7
    \item \textbf{Power delivery}: The Ultra 7 258V may maintain higher sustained clock speeds during the tests due to better thermal characteristics
\end{itemize}

\textbf{Conclusion}: Core count is not the determining factor for this workload. The newer Lunar Lake architecture (Ultra 7 258V) provides superior real-world performance despite being positioned as a low-power mobile processor. For reproducibility, we recommend documenting both hardware and complete software environment, as runtime performance can vary significantly across different CPU architectures.

\section{Conclusion}
This report presents a unified IEM pipeline with three core algorithms: Monte Carlo IC evaluation, MC-SAA greedy heuristic with IMRank screening, and pymoo-based binary evolutionary optimization. The formulation, pseudocode, and batch experiment protocol support reproducible comparison across provided and generated datasets. After filling the reserved result tables, the report can be directly used to analyze efficiency-quality trade-offs between heuristic and evolutionary strategies under consistent evaluation settings.

\bibliography{references} 
\bibliographystyle{IEEEtran} 
        


\end{document}
